%%
%% SIT330-770 Natural Language Processing
%% HD Task 2.3 - Research Paper
%% ACM Conference Proceedings Primary Article Template
%%

\documentclass[sigconf]{acmart}

\setcopyright{none}
\settopmatter{printacmref=false, printccs=false, printfolios=true}
\renewcommand\footnotetextcopyrightpermission[1]{}
\acmConference{SIT330-770 Natural Language Processing}{2025}{Deakin University}
\acmDOI{}
\acmISBN{}

\usepackage{booktabs}
\usepackage{graphicx}
\usepackage{microtype}

\title{Noise-Robust Sentiment Classification via\\
Noise-Augmented Fine-Tuning}

\author{Ali Alabdouli}
\affiliation{\institution{Deakin University}\country{Australia}}
\email{224535032}

\begin{document}

\begin{abstract}
Transformer-based sentiment classifiers trained on formal corpora
such as SST-2 degrade under informal text noise encountered in
real-world deployment.
While the existence of this degradation is acknowledged, the
relative severity of individual noise types and the effectiveness
of mitigation strategies remain understudied.
This paper makes two contributions.
First, we systematically quantify how six noise types (typographic
errors, abbreviations, slang, emoji insertion, missing punctuation,
and sarcasm) degrade sentiment classification performance at
incremental intensities $\alpha \in [0, 1]$.
Second, we propose and evaluate \emph{noise-augmented fine-tuning}
(NAFT), which exposes a pre-trained DistilBERT classifier to
synthetic noise during fine-tuning to improve robustness at
inference time.
Experiments on SST-2 show that typographic noise is the most
damaging noise type, reducing baseline macro-F1 from 0.91 to 0.33
at full intensity, and that NAFT recovers up to 0.29 F1 points
under this condition.
Emoji noise also shows substantial recovery under NAFT (+0.11 F1
at $\alpha=1.0$).
Surprisingly, the baseline model is already robust to slang,
abbreviations, and punctuation noise, while sarcasm as implemented
via prefix injection has minimal effect on either model.
These findings provide actionable guidance for practitioners
deploying sentiment systems in noisy environments.
\end{abstract}

\maketitle

%% ---------------------------------------------------------------
\section{Introduction}
%% ---------------------------------------------------------------

\paragraph{Background and Importance.}
Sentiment analysis is widely used to understand people's opinions
across domains including product reviews, customer feedback, social
media monitoring, and public opinion analysis~\cite{pang2008opinion}.
In real-world deployment settings, users frequently write
informally, using slang, abbreviations, emojis, sarcasm, spelling
mistakes, and missing punctuation~\cite{eisenstein2013bad}.
This informal register makes sentiment classification harder because
emotional meaning may not be recoverable from the surface form of
standard words alone.
A critical \emph{deployment gap} therefore exists: models are trained
in one context (clean, edited review text) but used in another
(unstructured, noisy social media), and the consequences of this
mismatch are not fully understood.

\paragraph{Current Approaches.}
Current sentiment analysis systems commonly employ transformer-based
models such as BERT~\cite{devlin2019bert} and
RoBERTa~\cite{liu2019roberta}, fine-tuned on labelled datasets.
Training corpora vary in register: SST-2~\cite{socher2013recursive}
consists of edited film review excerpts and represents a relatively
formal dataset, whereas IMDb~\cite{maas2011learning} comprises
raw user-submitted reviews containing slang and informal constructions.
Despite this variation, it remains unclear whether training register
confers meaningful noise robustness at inference time.

\paragraph{Identified Limitation.}
The relative severity of individual noise types has not been
systematically quantified in isolation, and no controlled study
has examined whether noise-augmented training directly mitigates
the observed degradation.
Without this evidence, practitioners cannot make informed decisions
about robustness strategies for informal deployment settings.

\paragraph{Research Question.}
\emph{Does noise-augmented fine-tuning reduce the performance
degradation caused by informal text noise in sentiment
classification, and does its effectiveness vary across noise types
and intensities?}

\paragraph{Contributions.}
(1) A controlled noise sensitivity benchmark across six noise types
at eleven intensity levels; (2) a noise-augmented fine-tuning method
(NAFT) that directly targets the identified limitation; and (3) an
empirical comparison demonstrating where NAFT recovers performance
and where it does not.

%% ---------------------------------------------------------------
\section{Related Work}
%% ---------------------------------------------------------------

\paragraph{Data Augmentation for Robustness.}
The principle that exposing a model to perturbed inputs during
training improves robustness to those perturbations at inference
time is well established in machine learning, predating
transformer models.
\citet{simard2003best} demonstrated this for convolutional networks
in computer vision, showing that elastic distortions during training
substantially improved generalisation to unseen image variations.
\citet{goodfellow2015explaining} formalised adversarial training,
arguing that training on worst-case perturbations acts as a
regulariser that improves robustness to input variation more
broadly.
NAFT applies this principle to the text domain with a
domain-specific perturbation taxonomy.

\paragraph{Text Data Augmentation.}
\citet{wei2019eda} introduced EDA (Easy Data Augmentation), which
applies four word-level operations: synonym replacement, random
insertion, random swap, and random deletion, as augmentation
strategies for text classification.
EDA showed that training with augmented data matched full-data
accuracy using only 50\% of the original training set.
NAFT is conceptually related to EDA in that both generate
perturbed training examples to improve robustness.
The key distinction is that EDA applies generic word-level
operations designed for data efficiency, while NAFT applies
noise types that directly mirror informal social media language
(typographic errors, slang, emoji insertion), targeting deployment
robustness rather than label efficiency.
EDA does not address emoji, typographic, or sarcasm perturbations,
which are the noise types most damaging in real-world deployment.

\paragraph{Noise Robustness in NLP.}
\citet{belinkov2018synthetic} showed that neural machine translation
systems are highly sensitive to character-level noise, and that
noise-aware training recovers much of the lost performance,
directly motivating NAFT for sentiment classification.
In sentiment analysis, \citet{barbieri2020tweeteval} showed
significant performance gaps between models evaluated on formal and
informal corpora, highlighting the deployment mismatch addressed
here.
Sarcasm and irony detection require pragmatic context beyond
token-level signals~\cite{van2018semeval}, a limitation confirmed
by our results.

%% ---------------------------------------------------------------
\section{Preliminaries and Problem Definition}
%% ---------------------------------------------------------------

Let $f: \mathcal{X} \to \{0, 1\}$ denote a binary sentiment
classifier mapping input text to a negative (0) or positive (1)
label.
Let $\mathcal{N}_k^\alpha: \mathcal{X} \to \mathcal{X}$ denote a
noise function for type $k$ at intensity $\alpha \in [0, 1]$,
where $\alpha$ represents the probability that each eligible token
position is perturbed.
We define six noise types $k \in \{\text{typo, abbrev, slang,
emoji, punct, sarcasm}\}$, each operationalised as follows.
\textbf{Typographic} noise ($\mathcal{N}_{\text{typo}}^\alpha$)
applies random character substitution, deletion, or duplication to
each alphabetic character with probability $\alpha$.
\textbf{Abbreviation} ($\mathcal{N}_{\text{abbrev}}^\alpha$) and
\textbf{slang} ($\mathcal{N}_{\text{slang}}^\alpha$) replace each
word present in a fixed lexicon with its informal equivalent with
probability $\alpha$.
\textbf{Emoji} ($\mathcal{N}_{\text{emoji}}^\alpha$) inserts a
sentiment-neutral emoji at each token boundary with probability
$\alpha$.
\textbf{Punctuation} ($\mathcal{N}_{\text{punct}}^\alpha$) removes
each punctuation character with probability $\alpha$.
\textbf{Sarcasm} ($\mathcal{N}_{\text{sarc}}^\alpha$) prepends a
pragmatic inversion frame to a fraction $\alpha$ of positive samples.

The degradation of classifier $f$ under noise type $k$ at intensity
$\alpha$ is measured as:
\begin{equation}
  \Delta F_1(f, k, \alpha) = F_1(f, \mathcal{D}) -
  F_1(f, \mathcal{N}_k^\alpha(\mathcal{D}))
\end{equation}
where $\mathcal{D}$ is the clean test set and $F_1$ is macro-F1.
The \emph{degradation threshold} $\alpha^*_k$ is the smallest
$\alpha$ at which $\Delta F_1 \geq 0.05$.

\paragraph{Noise Operationalisation Details.}
Each noise function is deterministic given a fixed random seed.
\textbf{Typographic} noise independently applies one of three
character-level operations to each alphabetic character with
probability $\alpha$: (i) substitution with a uniformly sampled
lowercase letter, (ii) deletion, or (iii) duplication of the
character.
\textbf{Abbreviation} noise uses a fixed lexicon of 30 common
English informal substitutions (e.g., \textit{you}$\to$\textit{u},
\textit{because}$\to$\textit{bc}, \textit{through}$\to$\textit{thru}).
Each word whose lowercase form appears in the lexicon is replaced
with probability $\alpha$; other words are unchanged.
\textbf{Slang} noise follows the same mechanism with a 16-entry
sentiment-relevant lexicon (e.g., \textit{good}$\to$\textit{lit},
\textit{bad}$\to$\textit{trash}, \textit{excellent}$\to$\textit{goated}).
\textbf{Emoji} noise inserts one of ten manually selected
sentiment-neutral emojis (face-with-tears-of-joy, eyes, hundred-points,
thinking-face, grinning-face-with-sweat, upside-down-face,
clapping-hands, sparkles, fire, skull) uniformly at random after
each token with probability $\alpha$.
\textbf{Punctuation removal} drops each character in
\{.,!?;:\} independently with probability $\alpha$.
\textbf{Sarcasm} prepends one of four fixed phrases
(\textit{``Oh sure,''} \textit{``Yeah right,''} \textit{``Totally,''} \textit{``Oh absolutely,''})
chosen uniformly at random, applied only to positive-labelled
samples with probability $\alpha$.

%% ---------------------------------------------------------------
\section{Proposed Method}
%% ---------------------------------------------------------------

We propose \textbf{Noise-Augmented Fine-Tuning (NAFT)}, which
addresses the identified limitation by exposing the model to
synthetic noise during fine-tuning, so that noisy token
distributions are in-distribution at inference time.

\paragraph{Method.}
Starting from a pre-trained DistilBERT model fine-tuned on SST-2,
we construct an augmented training set $\tilde{\mathcal{D}}$ by
pairing each original training example $(x, y) \in \mathcal{D}$
with noisy copies generated by applying all six noise types
simultaneously at each of three augmentation intensities
$\alpha \in \{0.2, 0.4, 0.6\}$.
The combined noise function applies all perturbations sequentially.
This expands the training set by a factor of four (one clean copy
plus three noisy copies per example).
The model is then fine-tuned on $\tilde{\mathcal{D}}$ for three
epochs using AdamW with learning rate $2 \times 10^{-5}$ and
weight decay $0.01$.

\paragraph{Justification.}
The core insight is that noise-augmented training reduces the
distributional shift between training and inference.
By explicitly including perturbed text during optimisation, the
model's attention heads are encouraged to attend to semantic
content rather than surface-level lexical patterns, making the
learned representations more robust to token-level perturbations.
This is analogous to data augmentation in computer vision, where
exposing models to transformed images improves generalisation to
unseen transformations~\cite{shorten2019survey}.

\paragraph{Limitations of the Approach.}
NAFT is not expected to recover performance for pragmatic noise
(sarcasm), because sarcasm operates at the discourse level rather
than the token level; token-level augmentation cannot teach the
model to invert sentiment labels based on conversational framing.
This motivates separate treatment of sarcasm as a distinct
sub-problem in future work.

%% ---------------------------------------------------------------
\section{Experimental Setup}
%% ---------------------------------------------------------------

\paragraph{Datasets.}
We use the SST-2 validation set (872 examples, balanced to 500 for
efficiency) as our evaluation set.
The SST-2 training set (67K examples, subsampled to 2,000 for
fine-tuning) provides the augmented training data.

\paragraph{Models.}
The \textbf{baseline} is
\texttt{distilbert-base-uncased-finetuned-sst-2-english}, a
DistilBERT model fine-tuned on SST-2.
The \textbf{NAFT model} is the same checkpoint further fine-tuned
on the noise-augmented training set described in Section 4.

\paragraph{Evaluation Protocol.}
Both models are evaluated on the same 500-example test set under
each of the six noise types at all eleven intensity levels
$\alpha \in \{0.0, 0.1, \ldots, 1.0\}$.
Macro-F1 and accuracy are reported at each $(\text{noise type},
\alpha)$ combination.
The clean baseline ($\alpha = 0$) is identical for both models
prior to fine-tuning; any change in clean performance after NAFT
fine-tuning is also reported to assess the clean-noisy trade-off.

\paragraph{Implementation.}
All experiments are implemented in Python using HuggingFace
Transformers~\cite{wolf2020transformers}.
Training uses AdamW with learning rate $2 \times 10^{-5}$, batch
size 16, and 3 epochs.
Results are averaged over a fixed random seed (42) for
reproducibility.

%% ---------------------------------------------------------------
\section{Results}
%% ---------------------------------------------------------------

Table~\ref{tab:results} reports macro-F1 for both models across
all eleven $\alpha$ levels for the two most sensitive noise types;
Table~\ref{tab:results2} covers the remaining four.\footnote{Code,
trained model checkpoints, and all noise-injected dataset variants
are available at:
\texttt{https://github.com/AliAlabdouli011/naft-sentiment}
(to be made public upon submission).}

\paragraph{Baseline Noise Sensitivity.}
Typographic noise is by far the most damaging noise type.
The baseline macro-F1 drops from 0.9119 at $\alpha=0.0$ to 0.7538
at $\alpha=0.1$ — a drop of 0.16 at the lowest non-zero intensity —
and collapses to 0.3333 at $\alpha \geq 0.9$, equivalent to
random guessing.
This rapid degradation occurs because character-level perturbations
produce tokens outside DistilBERT's WordPiece vocabulary,
fragmenting words into meaningless subword units and destroying
contextual representations.
The degradation threshold $\alpha^* = 0.1$ is the lowest of all
noise types, confirming typographic noise as the primary
vulnerability.

Emoji insertion causes moderate but steady degradation, with
baseline F1 declining from 0.9119 to 0.7621 at $\alpha=1.0$
($\Delta = -0.15$).
The model was not exposed to emoji tokens during SST-2 training,
so each inserted emoji disrupts attention over surrounding words.

Abbreviations, slang, and punctuation removal cause minimal
degradation.
Abbreviation F1 drops from 0.9119 to 0.8760 ($\Delta=-0.036$)
and slang from 0.9119 to 0.9000 ($\Delta=-0.012$) at full
intensity, likely because only a small fraction of tokens in any
sentence match the fixed lexicons, leaving sufficient context for
correct prediction.

Sarcasm prefix injection has negligible effect across all $\alpha$
levels (F1 remains $\approx0.91$), discussed in
Section~\ref{sec:analysis}.

\paragraph{NAFT Performance.}
NAFT provides the largest absolute improvements for typographic
noise, recovering up to $\Delta F_1 = +0.29$ at $\alpha=0.3$
(Table~\ref{tab:results}).
The improvement is consistent across all non-zero intensities,
confirming that character-level augmentation during training
significantly mitigates the tokenizer fragmentation problem.
Even at full intensity ($\alpha=1.0$), NAFT achieves F1=0.3838
versus baseline 0.3333 (+0.051).

Emoji noise also benefits substantially from NAFT, improving from
0.7621 to 0.8676 at $\alpha=1.0$ (+0.106), with gains growing
monotonically with intensity (Table~\ref{tab:results}).
This suggests the model learns to treat emoji tokens as ignorable
noise rather than semantic content.

For abbreviation, slang, and punctuation, NAFT produces small
negative deltas ($\leq -0.02$) at most levels.
Since the baseline is already robust to these types, augmentation
provides no benefit and introduces a marginal cost from the
clean-text trade-off.

\paragraph{Compound Noise.}
Table~\ref{tab:compound} reports results when all six noise types
are applied simultaneously.
The baseline degrades catastrophically, dropping from F1=0.9119 at
$\alpha=0.0$ to F1=0.4018 at $\alpha=1.0$.
NAFT, by contrast, maintains strong performance across the entire
range and reaches F1=0.9900 at $\alpha=1.0$ — substantially
\emph{exceeding} its clean baseline.
This remarkable result is explained by the training procedure:
NAFT was fine-tuned on exactly this type of combined noise via
\texttt{apply\_all\_noise}, making heavily perturbed compound
inputs effectively in-distribution at inference time.

\begin{table}[h]
  \centering\footnotesize\setlength{\tabcolsep}{4pt}
  \caption{Macro-F1 under compound noise (all 6 types simultaneously).}
  \label{tab:compound}
  \begin{tabular}{rccr}
    \toprule
    $\alpha$ & \textbf{Baseline} & \textbf{NAFT} & $\Delta$ \\
    \midrule
    0.0 & 0.9119 & 0.8959 & $-$0.016 \\
    0.1 & 0.7550 & 0.8120 & $+$0.057 \\
    0.3 & 0.4582 & 0.7164 & $+$0.258 \\
    0.5 & 0.3628 & 0.7323 & $+$0.370 \\
    0.7 & 0.3762 & 0.8093 & $+$0.433 \\
    1.0 & 0.4018 & 0.9900 & $+$0.588 \\
    \bottomrule
  \end{tabular}
\end{table}

\paragraph{Comparison with EDA.}
Table~\ref{tab:eda} compares NAFT against EDA~\cite{wei2019eda}
for typographic noise, the most discriminating condition.
EDA provides negligible improvement over the baseline (F1=0.43 vs
0.42 at $\alpha=0.3$), confirming that word-level operations
(synonym replacement, random insertion, swap, deletion) do not
address character-level corruption.
NAFT substantially outperforms both (F1=0.64 at $\alpha=0.3$,
$+$0.22 over EDA), demonstrating that domain-specific noise
augmentation is necessary for typographic robustness.
For abbreviation, slang, punctuation, and sarcasm, all three models
perform similarly since the baseline is already robust to these
types.

\begin{table}[h]
  \centering\footnotesize\setlength{\tabcolsep}{4pt}
  \caption{Macro-F1 comparison: Baseline vs EDA vs NAFT
           under typographic noise.}
  \label{tab:eda}
  \begin{tabular}{rcccc}
    \toprule
    $\alpha$ & \textbf{Base} & \textbf{EDA} & \textbf{NAFT} & \textbf{NAFT$-$EDA} \\
    \midrule
    0.0 & 0.9119 & 0.8959 & 0.8900 & $-$0.006 \\
    0.1 & 0.7329 & 0.7285 & 0.7779 & $+$0.049 \\
    0.2 & 0.5170 & 0.5697 & 0.6937 & $+$0.124 \\
    0.3 & 0.4207 & 0.4283 & 0.6406 & $+$0.212 \\
    0.5 & 0.3465 & 0.3490 & 0.5531 & $+$0.204 \\
    1.0 & 0.3333 & 0.3324 & 0.3917 & $+$0.059 \\
    \bottomrule
  \end{tabular}
\end{table}

\paragraph{Ablation over Training Intensities.}
Table~\ref{tab:ablation} reports an ablation over which
augmentation intensity $\alpha$ contributes most to NAFT's
robustness.
NAFT-0.2 and NAFT-0.4 perform comparably and best at low-to-moderate
noise ($\alpha \leq 0.5$), while NAFT-0.6 overfits to high-intensity
noise and underperforms at mild conditions.
The full NAFT model (trained at all three intensities) provides the
most consistent performance across the range, confirming that
multi-intensity augmentation is necessary: no single $\alpha$ alone
achieves the coverage of the combined schedule.

\begin{table}[h]
  \centering\footnotesize\setlength{\tabcolsep}{3pt}
  \caption{Ablation: macro-F1 under typographic noise for models
           trained at a single augmentation intensity vs all three.}
  \label{tab:ablation}
  \begin{tabular}{rccccc}
    \toprule
    $\alpha$ & \textbf{Base} & \textbf{N-0.2} & \textbf{N-0.4} &
               \textbf{N-0.6} & \textbf{N-all} \\
    \midrule
    0.0 & 0.9119 & 0.8980 & 0.8980 & 0.9079 & 0.8980 \\
    0.1 & 0.7329 & 0.7858 & 0.7926 & 0.7708 & 0.7996 \\
    0.3 & 0.4207 & 0.6049 & 0.5781 & 0.4995 & 0.6381 \\
    0.5 & 0.3465 & 0.5187 & 0.5312 & 0.4320 & 0.5602 \\
    0.7 & 0.3324 & 0.4447 & 0.4311 & 0.3597 & 0.4322 \\
    1.0 & 0.3333 & 0.4104 & 0.4098 & 0.3403 & 0.3737 \\
    \bottomrule
  \end{tabular}
\end{table}
NAFT reduces clean-text F1 from 0.9119 to 0.8998 ($\Delta=-0.012$),
a consistent but minimal drop observed across all noise types at
$\alpha=0$.
This represents mild catastrophic forgetting from fine-tuning on
noisy data, and constitutes a negligible cost relative to the
gains under typographic and emoji noise.

\begin{table}[h]
  \centering\footnotesize\setlength{\tabcolsep}{3pt}
  \caption{Macro-F1 across all $\alpha$ for typographic and emoji
           noise. $\Delta$ = NAFT $-$ Baseline.}
  \label{tab:results}
  \begin{tabular}{ccccr}
    \toprule
    \textbf{Noise} & $\alpha$ & \textbf{Base} & \textbf{NAFT} & $\Delta$ \\
    \midrule
    Typographic & 0.0 & 0.9119 & 0.8998 & $-$0.012 \\
                & 0.1 & 0.7538 & 0.7958 & $+$0.042 \\
                & 0.2 & 0.5524 & 0.7379 & $+$0.186 \\
                & 0.3 & 0.3821 & 0.6707 & $+$0.289 \\
                & 0.4 & 0.3773 & 0.6134 & $+$0.236 \\
                & 0.5 & 0.3456 & 0.5993 & $+$0.254 \\
                & 0.6 & 0.3359 & 0.5460 & $+$0.210 \\
                & 0.7 & 0.3324 & 0.4770 & $+$0.145 \\
                & 0.8 & 0.3378 & 0.4216 & $+$0.084 \\
                & 0.9 & 0.3333 & 0.4002 & $+$0.067 \\
                & 1.0 & 0.3333 & 0.3838 & $+$0.051 \\
    \addlinespace
    Emoji       & 0.0 & 0.9119 & 0.8998 & $-$0.012 \\
                & 0.1 & 0.9060 & 0.9099 & $+$0.004 \\
                & 0.2 & 0.8938 & 0.8938 & $+$0.000 \\
                & 0.3 & 0.8798 & 0.9058 & $+$0.026 \\
                & 0.4 & 0.8856 & 0.9119 & $+$0.026 \\
                & 0.5 & 0.8593 & 0.8918 & $+$0.033 \\
                & 0.6 & 0.8487 & 0.8838 & $+$0.035 \\
                & 0.7 & 0.8207 & 0.8918 & $+$0.071 \\
                & 0.8 & 0.8194 & 0.8777 & $+$0.058 \\
                & 0.9 & 0.7983 & 0.8715 & $+$0.073 \\
                & 1.0 & 0.7621 & 0.8676 & $+$0.106 \\
    \bottomrule
  \end{tabular}
\end{table}

\begin{table}[h]
  \centering\footnotesize\setlength{\tabcolsep}{3pt}
  \caption{Macro-F1 across all $\alpha$ for remaining noise types.}
  \label{tab:results2}
  \begin{tabular}{ccccr}
    \toprule
    \textbf{Noise} & $\alpha$ & \textbf{Base} & \textbf{NAFT} & $\Delta$ \\
    \midrule
    Abbrev. & 0.0 & 0.9119 & 0.8998 & $-$0.012 \\
            & 0.3 & 0.8920 & 0.8939 & $+$0.002 \\
            & 0.5 & 0.8820 & 0.8839 & $+$0.002 \\
            & 0.7 & 0.8800 & 0.8779 & $-$0.002 \\
            & 1.0 & 0.8760 & 0.8659 & $-$0.010 \\
    \addlinespace
    Slang   & 0.0 & 0.9119 & 0.8998 & $-$0.012 \\
            & 0.3 & 0.9140 & 0.9019 & $-$0.012 \\
            & 0.5 & 0.9020 & 0.8959 & $-$0.006 \\
            & 0.7 & 0.9080 & 0.8919 & $-$0.016 \\
            & 1.0 & 0.9000 & 0.8879 & $-$0.012 \\
    \addlinespace
    Punct.  & 0.0 & 0.9119 & 0.8998 & $-$0.012 \\
            & 0.3 & 0.9060 & 0.9059 & $-$0.000 \\
            & 0.5 & 0.9080 & 0.8999 & $-$0.008 \\
            & 0.7 & 0.9020 & 0.9020 & $+$0.000 \\
            & 1.0 & 0.8940 & 0.8960 & $+$0.002 \\
    \addlinespace
    Sarcasm & 0.0 & 0.9119 & 0.8998 & $-$0.012 \\
            & 0.3 & 0.9139 & 0.9078 & $-$0.006 \\
            & 0.5 & 0.9139 & 0.9078 & $-$0.006 \\
            & 0.7 & 0.9119 & 0.9257 & $+$0.014 \\
            & 1.0 & 0.9159 & 0.9297 & $+$0.014 \\
    \bottomrule
  \end{tabular}
\end{table}

%% ---------------------------------------------------------------
\section{Critical Analysis}
\label{sec:analysis}
%% ---------------------------------------------------------------

\paragraph{Error Analysis.}
The most informative failure mode is the typographic noise
collapse.
At $\alpha \geq 0.3$, baseline F1 stabilises near 0.33 — the
score of a degenerate classifier that always predicts the majority
class.
This suggests the model is not merely degraded but completely
non-functional: character-level perturbations fragment tokens so
severely that the input becomes uninterpretable to the model.
NAFT significantly delays this collapse, maintaining F1 above 0.60
up to $\alpha=0.5$, but ultimately faces the same ceiling at high
intensities because even augmented training cannot prepare the
model for arbitrarily corrupted vocabulary.

The sarcasm results require specific explanation.
Our implementation prepends a fixed phrase (e.g.,
\textit{``Oh sure,''}) to positive samples, but the baseline
model correctly classifies these inputs because it attends to the
underlying film-domain vocabulary (\textit{``entertaining''},
\textit{``strong performances''}) rather than the pragmatic
framing.
This reveals a limitation of our sarcasm operationalisation rather
than genuine robustness to sarcasm: real sarcasm involves subtle
contextual inversions that cannot be captured by prefix injection.
The NAFT model's slight improvement under sarcasm ($+$0.014 at
$\alpha=1.0$) is likely noise in the evaluation rather than a
meaningful signal.

\paragraph{Limitations.}
Three limitations of the present study warrant discussion.
First, the abbreviation and slang lexicons are small and fixed,
covering only a fraction of real informal vocabulary.
Results for these noise types may therefore underestimate
degradation in genuine social media text.
Second, noise types are applied independently at evaluation time,
whereas informal text in practice combines multiple noise types
simultaneously; joint noise conditions may produce non-additive
degradation patterns not captured here.
Third, the sarcasm operationalisation does not reflect genuine
discourse-level irony, and the robustness observed should not be
interpreted as sarcasm detection capability.

\paragraph{Generalisation.}
The experimental findings are specific to binary sentiment
classification on SST-2 with a DistilBERT backbone.
Whether NAFT gains transfer to multi-class sentiment, other
transformer architectures, or non-English informal text remains
an open empirical question.
The relative severity ordering (typographic $\gg$ emoji $>$
abbreviation $\approx$ slang $\approx$ punctuation) may also shift
on larger or more diverse lexicons.

\paragraph{Practical Implications.}
For practitioners deploying sentiment systems on informal text,
the results suggest three actionable conclusions.
First, typographic noise is the highest-priority robustness concern
and NAFT is an effective remedy at moderate intensities.
Second, emoji robustness is achievable through noise augmentation
at low cost.
Third, slang and abbreviation handling should focus on lexicon
expansion rather than model-level interventions, as the baseline
is already robust at the lexicon sizes tested.

%% ---------------------------------------------------------------
\section{Conclusion and Future Work}
%% ---------------------------------------------------------------

This paper investigated the noise sensitivity of a DistilBERT
sentiment classifier across six informal text noise types and
proposed noise-augmented fine-tuning (NAFT) as a principled remedy.
Typographic noise is the most damaging single noise type (F1
collapse from 0.91 to 0.33), and NAFT recovers up to 0.29 F1
points under this condition.
Under compound noise, NAFT's advantage is even more pronounced,
achieving F1=0.99 at full intensity versus the baseline's 0.40 —
a direct consequence of training on combined noise distributions.
Comparison with EDA confirms that word-level augmentation alone is
insufficient for character-level robustness, motivating
domain-specific noise taxonomies in future augmentation strategies.
The ablation study shows that multi-intensity training
($\alpha \in \{0.2, 0.4, 0.6\}$) outperforms any single intensity,
with low-intensity variants dominating at mild noise and the full
schedule providing the broadest coverage.

Future work should explore: (1) sarcasm-aware training using
irony-labelled datasets such as SemEval-2018 Task
3~\cite{van2018semeval}; (2) expanding the slang and abbreviation
lexicons to better reflect real social media vocabulary; and (3)
extending NAFT to multilingual informal text and other sentiment
model architectures.

%% ---------------------------------------------------------------
\bibliographystyle{ACM-Reference-Format}
\bibliography{references}

\end{document}
\end{document}
